% Generated draft: requires mathematical and source review.
\documentclass[11pt]{article}
\usepackage[margin=25mm]{geometry}
\usepackage{fontspec}
\setmainfont{Latin Modern Roman}
\usepackage{amsmath,amssymb,mathtools,unicode-math}
\setmathfont{Latin Modern Math}
\usepackage{enumitem,needspace,xurl,hyperref}
\setlength{\parindent}{0pt}
\setlength{\parskip}{5pt}
\setlength{\emergencystretch}{3em}
\newcommand{\N}{\mathbb{N}}
\newcommand{\Z}{\mathbb{Z}}
\newcommand{\Q}{\mathbb{Q}}
\newcommand{\R}{\mathbb{R}}
\newcommand{\C}{\mathbb{C}}
\newcommand{\F}{\mathbb{F}}
\newcommand{\A}{\mathbb{A}}
\newcommand{\PP}{\mathbb{P}}
\newcommand{\E}{\mathbb{E}}
\newcommand{\eps}{\varepsilon}
\newcommand{\dd}{\,\mathrm{d}}
\begin{document}
\section*{UnsolvedMath ID: 700764}
\textbf{Source problem code:} catalogue entry
\subsection{Definitions and mathematical statement}
A variety of groups \(\mathfrak V\) is the class of all groups satisfying a specified collection of group identities. For a positive integer \(r\), let \(F_r(\mathfrak V)\) denote the relatively free group of rank \(r\) in \(\mathfrak V\): it has distinguished generators \(x_1,\ldots,x_r\), and every assignment of these generators to elements of any group in \(\mathfrak V\) extends uniquely to a homomorphism. An embedding is an injective group homomorphism. A group is \(n\)-generated if it admits a generating set of size at most \(n\); countable includes finite groups.

The source calls \(\mathfrak V\) \emph{non-regular} if \(F_{n+1}(\mathfrak V)\) embeds in \(F_n(\mathfrak V)\) for some \(n\). Using the conventional interpretation that \(n\) is a positive integer, the question is as follows. Let \(\mathfrak V\) be a variety of groups and let \(n\geq 1\) satisfy
\[
F_{n+1}(\mathfrak V)\hookrightarrow F_n(\mathfrak V).
\]
Must it follow that
\[
\forall G\in\mathfrak V,\quad |G|\leq\aleph_0
\ \Longrightarrow\ 
\exists H\in\mathfrak V\;\exists h_1,\ldots,h_n\in H:
\quad H=\langle h_1,\ldots,h_n\rangle
\quad\text{and}\quad G\hookrightarrow H?
\]
Here the conclusion uses the same \(n\) as the assumed embedding of relatively free groups. The source does not explicitly specify the range of \(n\); the positive-integer convention is an interpretation requiring review.
\subsection{Short English statement}
If the relatively free group of rank \(n+1\) in a variety embeds in its relatively free group of rank \(n\), must every countable group in that variety embed in a group in the same variety generated by at most \(n\) elements? The same \(n\) is used throughout.
\subsection{Sources}
\begin{itemize}
\item Supplied local UnsolvedMath registry record, permanent ID \(700764\), exact source code \texttt{catalogue entry}, title \emph{Kourovka Notebook Problem 11.72}; upstream revision \texttt{d371ee917ac9647fbb06cd6fc529deec6f3dda82}. The record labels the problem open; this is supplied metadata, not an independently verified mathematical or literature assessment.
\item Attributed in the supplied record to A. Yu. Olshanskii, proposed in 1990, as Problem 11.72 in \emph{The Kourovka Notebook (Unsolved Problems in Group Theory)}, 11th issue (1990).
\item The supplied background also identifies the 21st issue edition (2026), arXiv:1401.0300v48. The supplied collection description names E. I. Khukhro and V. D. Mazurov as editors of the notebook covering Issues 1--21 (1965--2026). These bibliographic details are reproduced from the record and have not been independently checked.
\end{itemize}
\end{document}
